\documentclass[11pt]{article}
\usepackage[a4paper,margin=26mm]{geometry}
\usepackage[T1]{fontenc}
\usepackage{lmodern}
\usepackage{amsmath,amssymb,amsthm}
\usepackage{microtype}
\usepackage{booktabs}
\usepackage{xurl}
\usepackage{hyperref}
\hypersetup{colorlinks=true,linkcolor=blue,urlcolor=blue,citecolor=blue,
 pdftitle={On the rank of groups of equivariant permutations: a counterexample to an upper bound},
 pdfauthor={Horace Dong}}
\newtheorem{theorem}{Theorem}
\newtheorem{citedtheorem}{Theorem}
\renewcommand{\thecitedtheorem}{\Alph{citedtheorem}}
\newtheorem{lemma}{Lemma}
\newtheorem{proposition}{Proposition}
\newtheorem{corollary}{Corollary}
\theoremstyle{remark}
\newtheorem{remark}{Remark}
\newcommand{\repoversion}{1.23.0}
\setlength{\emergencystretch}{1.5em}

\newcommand{\Rank}{\operatorname{Rank}}
\newcommand{\Aut}{\operatorname{Aut}}
\newcommand{\Conj}{\operatorname{Conj}}
\newcommand{\Sym}{\operatorname{Sym}}
\newcommand{\lmin}{\lambda^{\min}}

\title{On the rank of groups of equivariant permutations:\\
a counterexample to an upper bound}
\author{Horace Dong\\ \small Independent researcher\\
\small\texttt{maxwelldhx@gmail.com}\\
\small ORCID: \href{https://orcid.org/0009-0002-7554-7548}{0009-0002-7554-7548}}
\date{19 November 2026}
\begin{document}
\maketitle

\begin{abstract}
For a finite group $G$ acting on a finite set $X$, Lara-G\'omez,
Romero-Rosales and Ruiz-Medina gave an upper bound for the rank of the group
$\Aut_G(X)$ of $G$-equivariant permutations of $X$, refined it by grouping
generators of coprime orders into ``packets'', and stated that the refined
bound should be the exact value of the rank. We point out that the first
bound, Theorem~4 of the published version, fails for the group of order two
acting regularly on two points, where it gives $0$ while the rank is $1$. We
give a five-point example in which the rank is $2$ while the refined
expression is $3$ (preprint version) or $4$ (published version); the authors'
own $21$-point example already shows that the refined expression is not
always exact, and ours is a simpler instance. Finally, for every $n\ge3$ we
give a faithful action of a finite cyclic group for which the rank is $2$
while the two versions of the refined expression equal $n$ and $n+1$ for
every admissible choice of generators, so the difference is unbounded. These
three results are formally verified in Lean~4. The proofs were found and
formalized by AI agents working under the author's direction; their role is
described at the end of the paper.
\end{abstract}

\section{Introduction}

Let a group $G$ act on a set $X$. The $G$-equivariant maps $f\colon X\to X$,
those with $f(g\cdot x)=g\cdot f(x)$ for all $g\in G$ and $x\in X$, form a
monoid under composition, and its group of units $\Aut_G(X)$ consists of the
$G$-equivariant permutations of $X$. Lara-G\'omez, Romero-Rosales and
Ruiz-Medina~\cite{lrr} study generating sets of this monoid and of its group of
units; an earlier version of their paper is the preprint~\cite{rlr}. For a
finite $G$-set $X$ they construct an explicit generating set of $\Aut_G(X)$
from the orbit structure of the action and count it, which gives an upper
bound for the rank of $\Aut_G(X)$ (Theorem~4 of~\cite{lrr}, Theorem~5
of~\cite{rlr}). They then group some of these generators into ``packets'' of
pairwise commuting elements of coprime orders, which gives a second, packet
bound (Corollary~6 of~\cite{lrr}, Corollary~19 of~\cite{rlr}), and in the
Open problem that closes the relevant subsection they write that the packet
bound ``should be the exact value of the rank''.

We make three observations; the definitions are recalled in
Section~\ref{sec:defs}.
\begin{itemize}
\item The printed upper bound of Theorem~4 of~\cite{lrr} (and the right-hand
side of Theorem~5 of~\cite{rlr}) fails for the cyclic group of order two
acting regularly on two points: the bound is $0$ and the rank is $1$
(Theorem~\ref{thm:two}). The same happens for every transitive action $G/H$
with $N_G(H)/H$ nontrivial (Remark~\ref{rem:transitive}, a written argument
only), and the bound becomes valid if the term $-\delta_G(X)$ is removed
(Remark~\ref{rem:repair}, also a written argument, which only recounts the
authors' own generating set).
\item The packet expression is not always the exact rank. This is already
shown by the authors' own $21$-point example (Example~1 of~\cite{lrr}), whose
printed generators give rank at most $4$ while the expression is $6$ ($5$ in
the version of~\cite{rlr}); see Section~\ref{sec:five}. We give a simpler
five-point example in which the rank is $2$ and the two versions of the
expression are $3$ and~$4$ (Theorem~\ref{thm:five}). The rank of the group
involved, $C_2\times S_3$, also follows from a theorem of Lu~\cite{lu}.
\item The difference between the packet expression and the rank is
unbounded: for every $n\ge3$ there is a faithful action of a finite cyclic
group with rank $2$ for which the two versions of the expression equal $n$ and
$n+1$, for every admissible choice of generating sets of minimal cardinality
(Theorem~\ref{thm:gap}).
\end{itemize}
To our knowledge, what is new here is the failure of Theorem~4
of~\cite{lrr} and the unbounded gap. In all three examples the packet
bounds~\eqref{eq:packet} hold. For the five-point example and the family both
versions of the expected equality fail; for the two-point example the journal
version fails ($2\ne1$), and so does Theorem~4 of~\cite{lrr}. We found no correction of~\cite{lrr}
or~\cite{rlr} in the literature as of 19 November 2026. Theorems~\ref{thm:two},
\ref{thm:five} and~\ref{thm:gap} are formally verified in Lean~4
(Section~\ref{sec:lean}); the remarks marked ``written argument'' are not.

\section{Definitions and the statements of~\texorpdfstring{\cite{lrr}}{[2]} and~\texorpdfstring{\cite{rlr}}{[5]}}
\label{sec:defs}

We follow the notation of~\cite{lrr}. Let $X$ be a finite $G$-set. For $x\in
X$ let $G_x$ be its stabilizer. Let $\Conj_G(X)$ be the set of conjugacy
classes $[H]$ of subgroups of $G$ that occur as stabilizers, and write
$r=|\Conj_G(X)|$. The \emph{box} of $[H]$ is
$B_{[H]}=\{x\in X: [G_x]=[H]\}$; it is a union of orbits. Put
\begin{align*}
\kappa_G(X)&=\#\{[H]\in\Conj_G(X): B_{[H]}\text{ consists of one orbit}\},\\
\omega_G(X)&=\#\{[H]\in\Conj_G(X): N_G(H)/H \text{ is trivial}\},\\
R_G(X)&=\textstyle\sum_{[H]\in\Conj_G(X)}\Rank\bigl(N_G(H)/H\bigr),
\end{align*}
where $\Rank$ is the least size of a generating set, with the convention
of~\cite{lrr,rlr} that the trivial group has rank $1$ (the notation $R_G(X)$
is ours). Further, $\delta_G(X)=0$ if $\omega_G(X)=R_G(X)$ and
$\delta_G(X)=1$ otherwise, and $\Lambda_G(X)=1$ if either every box consists of
one orbit and every $N_G(H)/H$ is trivial, or some box contains at least three
orbits; otherwise $\Lambda_G(X)=0$.

\begin{citedtheorem}[{\cite[Theorem~4]{lrr}; the right-hand side of
\cite[Theorem~5]{rlr}}]\label{thm:lrr}
Let $X$ be a finite $G$-set. Then
\[
\Rank\bigl(\Aut_G(X)\bigr)\le r-\kappa_G(X)+\Lambda_G(X)-\omega_G(X)
-\delta_G(X)+R_G(X).
\]
\end{citedtheorem}

The hypothesis is only that $X$ is a finite $G$-set. The displayed definition
of $\delta_G(X)$ is the same in both versions; the sentence explaining it says
that $\delta_G(X)$ equals zero when $\omega_G(X)=R_G(X)$ in~\cite[p.~8]{lrr}
and that it equals one in that case in~\cite[p.~6]{rlr}
(Remark~\ref{rem:delta}).

\paragraph{Packets.} For each box choose a generating set of minimal
cardinality, that is, with $\Rank(N_G(H)/H)$ elements, of $N_G(H)/H$ (the
authors' choice~\cite[p.~8]{lrr}, \cite[p.~8]{rlr}), with representatives
$n_1,\dots,n_k\in N_G(H)$, and a
representative orbit $Gx$ in the box with $G_x=H$. Each $n_t$ gives the
\emph{base generator} that maps $g\cdot x\mapsto gn_t\cdot x$ on $Gx$ and
fixes all other points; its order equals the order of $n_tH$ in $N_G(H)/H$.
Let $\mathcal B$ be the set of base generators of order greater than $1$. A
\emph{packet} is a subset of $\mathcal B$ whose elements have pairwise coprime
orders and come from pairwise different boxes, and $\lmin_G(X)$ is the least
number of packets in a partition of~$\mathcal B$. The authors note that
$\lmin_G(X)$ may depend on the choice of these generating
sets~\cite[p.~19]{lrr}. Throughout, ``admissible choice'' means such a choice;
generating sets that are only minimal under inclusion are not allowed. Write $L_G(X)=r-\kappa_G(X)+\Lambda_G(X)$. The packet
bounds are
\begin{equation}\label{eq:packet}
\begin{aligned}
\Rank\bigl(\Aut_G(X)\bigr)&\le L_G(X)+\lmin_G(X)&&\text{\cite[Corollary~19]{rlr}},\\
\Rank\bigl(\Aut_G(X)\bigr)&\le L_G(X)+\lmin_G(X)+1&&\text{\cite[Corollary~6]{lrr}}.
\end{aligned}
\end{equation}

\paragraph{The expectation of equality.} After Example~20 of~\cite{rlr}
(\S2.1, p.~19) and after Example~1 of~\cite{lrr} (\S3.1, p.~20), the Open
problem states that the respective bound in~\eqref{eq:packet} should be the
exact value of the rank of $\Aut_G(X)$, while noting that $\lmin_G(X)$ has no
explicit expression yet. We call this the authors' stated expectation of
equality: for every finite $G$-set $X$,
\[
\Rank\bigl(\Aut_G(X)\bigr)=L_G(X)+\lmin_G(X)\ \text{in \cite{rlr}},\qquad
\Rank\bigl(\Aut_G(X)\bigr)=L_G(X)+\lmin_G(X)+1\ \text{in \cite{lrr}}.
\]
The Introduction of~\cite{lrr} (p.~2) also says ``We conjecture that the upper
bound is sharp''. If ``sharp'' means only that the bound is attained for some
$G$-sets, our examples do not contradict that sentence; they concern the
stronger reading in the Open problem.

\section{The bound of Theorem~\ref{thm:lrr} fails}\label{sec:two}

\begin{theorem}\label{thm:two}
Let $G=C_2$ act regularly on a set $X$ of two points. Then $r=\kappa_G(X)=1$,
$\Lambda_G(X)=\omega_G(X)=0$, $R_G(X)=\delta_G(X)=\lmin_G(X)=1$, the
right-hand side of Theorem~\ref{thm:lrr} equals $0$, and
$\Rank(\Aut_G(X))=1$. So Theorem~\ref{thm:lrr} fails for this $G$-set. The
two packet bounds~\eqref{eq:packet} equal $1$ and $2$ and hold.
\end{theorem}

\begin{proof}
Both points lie in one orbit and have trivial stabilizer, so there is one box,
consisting of one orbit: $r=\kappa_G(X)=1$. Its quotient is
$N_G(1)/1=C_2$, of rank $1$, so $\omega_G(X)=0$, $R_G(X)=1$ and
$\delta_G(X)=1$. Neither condition for $\Lambda_G(X)=1$ holds, so
$\Lambda_G(X)=0$, and the right-hand side is $1-1+0-0-1+1=0$. The two
permutations of $X$, the identity and the swap, both commute with the action,
so $\Aut_G(X)\cong C_2$ has rank $1$. The only generating set of $C_2$ of minimal
cardinality is its nonidentity element, so $\mathcal B$ has one element and
$\lmin_G(X)=1$; hence $L_G(X)+\lmin_G(X)=1$ and $L_G(X)+\lmin_G(X)+1=2$.
\end{proof}

In the counting that proves Theorem~\ref{thm:lrr}~\cite[p.~16]{lrr}, the term
$-\delta_G(X)$ accounts for merging one base generator with the product $T$ of
the orbit cycles $t_i$. When every box has one orbit, $T$ is the identity
(the authors note on p.~10 that $s_i$ and $t_i$ can then be omitted), and the merge
does not save a generator; Theorem~\ref{thm:two} is the smallest such case.

\begin{remark}[Transitive actions; written argument, not formalized]
\label{rem:transitive}
Let $X=G/H$ be a finite transitive $G$-set with $W=N_G(H)/H$ nontrivial. An
equivariant permutation $f$ is determined by $f(H)=nH$, and $f$ and $f^{-1}$
preserve stabilizers, so $n\in N_G(H)$; conversely each $n\in N_G(H)$ gives
$gH\mapsto gnH$. Hence $\Aut_G(X)$ is isomorphic to the opposite group of $W$,
which is isomorphic to $W$ by inversion, and $\Rank(\Aut_G(X))=\Rank(W)$. Here
$r=\kappa_G(X)=1$, $\Lambda_G(X)=\omega_G(X)=0$, $R_G(X)=\Rank(W)$ and
$\delta_G(X)=1$, so the right-hand side of Theorem~\ref{thm:lrr} is
$\Rank(W)-1$. In particular the bound fails for the regular action of every
nontrivial finite group.
\end{remark}

\begin{remark}[A corrected bound; written argument, not formalized]
\label{rem:repair}
Removing $-\delta_G(X)$ gives a valid bound,
\[
\Rank\bigl(\Aut_G(X)\bigr)\le r-\kappa_G(X)+\Lambda_G(X)-\omega_G(X)+R_G(X),
\]
with equality in Remark~\ref{rem:transitive}. This is a recount of the
generating set of~\cite[Theorems~5 and~6]{lrr}: the $R_G(X)-\omega_G(X)$
nontrivial base generators, one transposition $s_i$ of two orbits in each of
the $r-\kappa_G(X)$ boxes with at least two orbits, and, if some box has at
least three orbits, one element $T$ that cyclically permutes the orbits of
every such box. Conjugates of $s_i$ by powers of $T$ generate all
permutations of the orbits of box $i$, and conjugating the base generators by
these gives the base generators on every orbit. If every box has one orbit and every
$N_G(H)/H$ is trivial, then $\Aut_G(X)$ is trivial, of rank $1$ by
convention, and the right-hand side is $\Lambda_G(X)=1$. We do not claim that
this bound is new or best possible.
\end{remark}

\begin{remark}[The explanatory sentence in~\cite{rlr}; written argument, not
formalized]\label{rem:delta}
If $\delta_G(X)$ is read as in the sentence of~\cite[p.~6]{rlr}, so that
$\delta_G(X)=1$ when $\omega_G(X)=R_G(X)$, take instead the trivial group
acting on two points. There is one box with two orbits and trivial quotient:
$r=1$, $\kappa_G(X)=\Lambda_G(X)=0$ and $\omega_G(X)=R_G(X)=1$. The bound is
then $1-0+0-1-1+1=0$, while $\Aut_G(X)=\Sym(X)$ has rank $1$. With the
displayed definition, $\delta_G(X)=0$ and the bound is $1$, which is correct.
So the right-hand side of Theorem~5 of~\cite{rlr} fails under either reading.
\end{remark}

\section{The packet expression is not always exact}\label{sec:five}

The authors' Example~1 of~\cite{lrr} (Example~20 of~\cite{rlr}) is a $G$-set
with $21$ points, two boxes of three orbits each, and quotients
$C_2\times C_2$ and $C_3$. Its parameters are $r=2$, $\kappa_G(X)=0$ and
$\Lambda_G(X)=1$; the two base generators of order $2$ must lie in different
packets and the base generator of order $3$ can join one of them, so
$\lmin_G(X)=2$ and the two versions of the packet expression equal $5$
and~$6$. The authors list four generators of $\Aut_G(X)$ and conclude that its
rank is at most $4$~\cite[p.~20]{lrr}. Read together with the Open problem
that follows it, this example already shows that the expected equality does
not always hold. (As a check outside the Lean formalization, we realized this
$G$-set with $G=C_2^2\times C_3$, each box consisting of three regular orbits
of the corresponding quotient, and found by computer that the four printed
elements, composed from right to left as in~\cite[Lemma~3]{lrr}, generate a
group of order $(4^3\cdot3!)(3^3\cdot3!)=62{,}208$, the
order of $\Aut_G(X)$ stated in~\cite{lrr}.) The following example is smaller
and can be checked by hand.

\begin{theorem}\label{thm:five}
Let $G=C_2$ act on $X=\{1,2,3,4,5\}$, its generator acting as the
transposition $(1\,2)$. Then $r=2$, $\kappa_G(X)=\Lambda_G(X)=\omega_G(X)=1$,
$R_G(X)=2$, $\delta_G(X)=\lmin_G(X)=1$, the right-hand side of
Theorem~\ref{thm:lrr} is $2$, and $\Rank(\Aut_G(X))=2$,
while $L_G(X)+\lmin_G(X)=3$ and $L_G(X)+\lmin_G(X)+1=4$.
\end{theorem}

\begin{proof}
The points $1,2$ form one orbit with trivial stabilizer, and $3,4,5$ are three
fixed points with stabilizer $G$. Since $G$ is abelian, $\Conj_G(X)$ has the
two elements $[1]$ and $[G]$, the first box has one orbit and the second
has three, so $r=2$, $\kappa_G(X)=1$ and $\Lambda_G(X)=1$. The quotients are
$C_2$ and the trivial group, each of rank $1$, so $\omega_G(X)=1$,
$R_G(X)=2$ and $\delta_G(X)=1$, and the right-hand side of
Theorem~\ref{thm:lrr} is $2-1+1-1-1+2=2$. Each quotient has a unique generating
set of minimal cardinality (for the trivial group, $\{1\}$, by the rank-$1$
convention); the base generator from the trivial quotient is the identity
and is discarded, and the one from $C_2$ is $(1\,2)$, of order $2$. So
$\lmin_G(X)=1$, and $L_G(X)+\lmin_G(X)=3$.

A permutation commutes with $(1\,2)$ if and only if it preserves $\{1,2\}$
and $\{3,4,5\}$, so $\Aut_G(X)=\langle(1\,2)\rangle\times\Sym\{3,4,5\}\cong
C_2\times S_3$. With $x=(1\,2)(3\,4\,5)$ and $y=(3\,4)$ we have
$x^3=(1\,2)$ and $x^4=(3\,4\,5)$, and $(3\,4\,5)$, $(3\,4)$ generate
$\Sym\{3,4,5\}$; so $x,y$ generate $\Aut_G(X)$. The group is not abelian,
hence not cyclic, and its rank is $2$.
\end{proof}

Here Theorem~\ref{thm:lrr} holds with equality and both packet bounds hold.
That $C_2\times S_3$ has rank $2$ is also a special case of Lu's theorem on
generalised wreath products of symmetric groups~\cite[Theorem~1.1 and
Lemma~3.10]{lu}; what Theorem~\ref{thm:five} adds is the computation of the
parameters of~\cite{lrr} for this $G$-set.

\paragraph{Two small points in~\cite{rlr}.} After Corollary~19 of~\cite{rlr}
(p.~18) the packet bound is compared with the earlier one using the inequality
$\lmin_G(X)\le R_G(X)-\omega_G(X)-\delta_G(X)$; for the $G$-set of
Theorem~\ref{thm:two} this reads $1\le0$. The journal version does not contain
this inequality. In Example~20 of~\cite{rlr} the set displayed before the
conclusion that the rank is at most $4$ has six elements; the journal version
lists four.

\section{An unbounded gap}\label{sec:gap}

\begin{theorem}\label{thm:gap}
For every $n\ge3$ there are a finite cyclic group $G$ and a finite set $X$ on
which $G$ acts faithfully such that $r=\kappa_G(X)=n$, $\Lambda_G(X)=0$,
$\Rank(\Aut_G(X))=2$, and, for every choice of generating sets of minimal
cardinality of the quotients $N_G(H)/H$, $\lmin_G(X)=n$. Hence the two versions of the packet
expression equal $n$ and $n+1$, and exceed the rank by $n-2$ and $n-1$.
\end{theorem}

\begin{proof}
Assign to each ordered pair $(i,j)$ of distinct elements of
$\{1,\dots,n\}$ a prime $p_{ij}$, all $n(n-1)$ primes distinct, and put
$p_{ii}=1$. Let
\[
A_i=\prod_{j}p_{ij},\qquad B_i=\prod_{j}p_{ji},\qquad q_i=A_iB_i,\qquad
M=\prod_iA_i=\prod_iB_i .
\]
For each $i$, $A_i$ and $B_i$ are coprime, since no label has $i$ as both its
first and its second index (except $p_{ii}=1$). The $A_i$ are pairwise coprime
and so are the $B_i$, so $A_i\mid M$ and $B_i\mid M$ give $q_i\mid M$. Each
$q_i>1$. For $i\ne j$ choose $k\notin\{i,j\}$ (here $n\ge3$ is used): $p_{ik}$
divides $q_i$ but not $q_j$, so the $q_i$ are distinct; and $p_{ij}$ divides
both $q_i$ and $q_j$.

Let $G=\mathbb Z/M$ act on $X=\bigsqcup_{i}\mathbb Z/q_i$ by
$g\cdot(i,t)=(i,t+g)$. If $g$ acts trivially, then $g\equiv0$ modulo every
$q_i$, hence modulo every $A_i$, hence modulo $M$ by the Chinese remainder
theorem; so the action is faithful. The reduction $\mathbb Z/M\to\mathbb Z/q_i$
is onto, so each component is one orbit, and the stabilizer of each of its
points is the kernel $K_i$ of the reduction. These kernels have the distinct
indices $q_i$, so they are distinct, and as $G$ is abelian each is its own
conjugacy class with normalizer $G$ and quotient $G/K_i\cong\mathbb Z/q_i$.
Thus each box is one component, consisting of one orbit: $r=\kappa_G(X)=n$.
No quotient is trivial and no box has three orbits, so $\Lambda_G(X)=0$.

A generating set of minimal cardinality of the nontrivial cyclic group
$\mathbb Z/q_i$ is a single generator, of order $q_i$. So for every choice $\mathcal B$ has
exactly one element from each box, of order $q_i$; any two of these orders
share the prime $p_{ij}$, so every packet has one element and
$\lmin_G(X)=n$.

An equivariant permutation preserves stabilizers and therefore maps each
component to itself. On $\mathbb Z/q_i$, equivariance gives
$f(t)=t+f(0)$, and conversely independent translations of the components are
equivariant. Hence, using the Chinese remainder theorem twice,
\[
\Aut_G(X)\cong\prod_i\mathbb Z/q_i\cong\prod_i\bigl(\mathbb Z/A_i\times
\mathbb Z/B_i\bigr)\cong\Bigl(\prod_i\mathbb Z/A_i\Bigr)\times
\Bigl(\prod_i\mathbb Z/B_i\Bigr)\cong\mathbb Z/M\times\mathbb Z/M .
\]
This group is generated by two elements and, since $M>1$, is not cyclic (every
element has order dividing $M<M^2$). So its rank is $2$.
\end{proof}

The cardinality condition matters here: $\{a_i,b_i\}$ with $a_i$ of order
$A_i$ and $b_i$ of order $B_i$ is an inclusion-minimal generating set of
$\mathbb Z/q_i$, and with these sets the packets $\{a_1,\dots,a_n\}$ and
$\{b_1,\dots,b_n\}$ would be allowed, giving $2$ instead of~$n$.

\begin{remark}[A smaller variant; written argument, not formalized]
One prime $p_{ij}=p_{ji}$ per unordered pair also works, with
$\Aut_G(X)\cong C_M\times C_M$, $M=\prod_{i<j}p_{ij}$; for $n=3$ and the primes
$2,3,5$ this gives orbits of lengths $6,10,15$ ($|X|=31$, $|G|=30$). Lean uses
the directed version.
\end{remark}


\section{Formal verification}\label{sec:lean}

Theorems~\ref{thm:two}, \ref{thm:five} and~\ref{thm:gap} are formally verified
in Lean~4 (version 4.34.1) with Mathlib~\cite{mathlib} (commit
\texttt{d13f23b7}). The $15$ files are in version~\repoversion{} of the
repository~\cite{repo}, directory \texttt{lean/Research/TaskJ/Rank/}: five files for
Theorems~\ref{thm:two} and~\ref{thm:five}, and ten, in a subdirectory
\texttt{Extension}, for Theorem~\ref{thm:gap}.

\paragraph{The statements.}
The two statement files \texttt{Challenge.lean} contain only definitions
and statements. There are seven frozen targets: \texttt{TwoPointParameters},
\texttt{FivePointParameters}, \texttt{TwoPointRank}, \texttt{FivePointRank},
\texttt{TwoPointCounterexample} and \texttt{FivePointCounterexample}, in the
namespace \texttt{TaskJ.\allowbreak Rank.\allowbreak Challenge}, and
\texttt{TaskJ.\allowbreak Rank.\allowbreak Extension.\allowbreak UnboundedGap}.
In ASCII notation, with implicit arguments omitted and some definitions
abridged:
{\footnotesize
\begin{verbatim}
-- Theorems 1 and 2: C2 = ZMod 2 acts on Fin n through an involution t
def twoSwap : Equiv.Perm (Fin 2) := Equiv.swap 0 1  -- fiveSwap: the same in Fin 5
def act (t : Equiv.Perm (Fin n)) (g : C2) (x : Fin n) : Fin n :=
  if g = 0 then x else t x
def equivariantGroup (t : Equiv.Perm (Fin n)) : Subgroup (Equiv.Perm (Fin n)) :=
  Subgroup.centralizer ({t} : Set _)
def firstBound t : Int :=
  (r t : Int) - kappa t + capitalLambda t - omega t - deltaDisplayed t + rankSum t
def preprintPacketBound t : Int :=
  (r t : Int) - kappa t + capitalLambda t + packetNumber t
def journalPacketBound t : Int := preprintPacketBound t + 1
def TwoPointParameters : Prop :=
  r twoSwap = 1 /\ kappa twoSwap = 1 /\ capitalLambda twoSwap = 0 /\
  omega twoSwap = 0 /\ rankSum twoSwap = 1 /\ deltaDisplayed twoSwap = 1 /\
  packetNumber twoSwap = 1 /\ firstBound twoSwap = 0 /\
  preprintPacketBound twoSwap = 1 /\ journalPacketBound twoSwap = 2
-- FivePointParameters: the values 2,1,1,1,2,1,1,2,3,4 for fiveSwap
def TwoPointRank : Prop := Group.rank (equivariantGroup twoSwap) = 1
def FivePointRank : Prop := Group.rank (equivariantGroup fiveSwap) = 2
def TwoPointCounterexample : Prop :=
  ActionLaws twoSwap /\ EquivarianceMeaning twoSwap /\ BasesValid twoSwap /\
  TwoPointParameters /\ TwoPointRank /\
  Not ((Group.rank (equivariantGroup twoSwap) : Int) <= firstBound twoSwap)
-- FivePointCounterexample: the analogous conjunction for fiveSwap, ending with
--   rank < preprintPacketBound fiveSwap /\ rank < journalPacketBound fiveSwap

-- Theorem 3: Z/M acts on the disjoint union of the Z/(width i)
structure Family (n : Nat) where
  modulus : PNat
  width : Fin n -> PNat
  divides : forall i, (width i : Nat) | (modulus : Nat)
def automorphisms (D : Family n) : Subgroup (Equiv.Perm D.Point) :=
  Subgroup.centralizer (Set.range D.action)
abbrev Weyl (D : Family n) (i : Fin n) := D.Actor / D.boxGroup i
def Valid (D : Family n) (B : D.Bases) : Prop :=
  forall i, Subgroup.closure (B i : Set (D.Weyl i)) = top /\
    (B i).card = max 1 (Group.rank (D.Weyl i))
def Fits (D : Family n) (B : D.Bases) (k : Nat) : Prop :=
  exists f : D.Item B -> Fin k, forall x y, x != y -> f x = f y ->
    x.1 != y.1 /\ Nat.Coprime (orderOf x.2.val) (orderOf y.2.val)
def packets (D : Family n) (B : D.Bases) : Nat := sInf {k : Nat | D.Fits B k}
-- preprint B := r - kappa + capitalLambda + packets B (in Int); journal B := preprint B + 1
def UnboundedGap : Prop :=
  forall n : Nat, 3 <= n -> exists D : Family n,
    (forall i, 1 < (D.width i : Nat)) /\ Function.Injective D.width /\
    Function.Injective D.action /\ IsCyclic D.Actor /\
    (forall x : D.Point, D.stabilizer x = (D.projection x.1).ker) /\
    (forall x : D.Point, D.orbit x = {y | y.1 = x.1}) /\
    Function.Injective D.boxGroup /\
    (forall i, Nonempty (MulEquiv (D.Weyl i) (D.Fiber i))) /\
    D.r = n /\ D.kappa = n /\ D.capitalLambda = 0 /\
    Group.rank D.automorphisms = 2 /\ (exists B : D.Bases, D.Valid B) /\
    forall B : D.Bases, D.Valid B -> D.packets B = n /\
      D.preprint B = (n : Int) /\ D.journal B = (n : Int) + 1
\end{verbatim}}
\noindent (In the files, \texttt{/} in \texttt{Weyl} is the quotient group,
\texttt{top} is the whole group and \texttt{!=} is $\ne$.) The equivariant
group is the centralizer of the action in the full permutation group (for
Theorems~\ref{thm:two} and~\ref{thm:five} the files also prove that this is
the set of permutations $p$ with $p(g\cdot x)=g\cdot p(x)$). Stabilizers,
orbits, boxes, the quotients $N_G(H)/H=G/H$ ($G$ is abelian), $r$,
$\kappa_G(X)$, $\omega_G(X)$, $R_G(X)$, $\Lambda_G(X)$, $\delta_G(X)$ (in its
displayed form) and $\lmin_G(X)$ are computed from the action as in
Section~\ref{sec:defs} (for Theorem~\ref{thm:gap} only $r$, $\kappa_G(X)$,
$\Lambda_G(X)$ and $\lmin_G(X)$ occur), and the expressions are evaluated in
the integers.
The rank of a quotient uses nonempty generating sets, matching the convention
that the trivial group has rank~$1$; Mathlib's \texttt{Group.rank} gives the
trivial group rank $0$, which does not matter as the equivariant groups are
nontrivial. For Theorems~\ref{thm:two} and~\ref{thm:five} the files fix one
generating set per quotient (\texttt{selectedBasis}) and prove that it
has minimal cardinality (\texttt{BasesValid}); this loses nothing, since the
quotients that occur, $C_2$ and the trivial group, each have a unique nonempty
generating set of minimal cardinality. \texttt{UnboundedGap} quantifies over every valid choice and
asserts that one exists; packets are encoded as colour classes. The Lean
statements give the parameter values, the inequalities and the ranks; the
isomorphism types named in the proofs ($C_2\times S_3$ and
$\mathbb Z/M\times\mathbb Z/M$) are written arguments and not part of them.

\paragraph{Review and freezing.}
The statement file for Theorems~\ref{thm:two} and~\ref{thm:five} was reviewed
against both versions of the source by a Codex agent that had not written it,
and frozen, by recording its SHA-256 hash, on 18 November 2026 at 00:07~UTC;
the statement file for Theorem~\ref{thm:gap} was reviewed by the same agent and
frozen at 00:22~UTC. The same agent had also checked the reading of the source
(see the last section). The development logs record no compilation of a
proof file before the corresponding freeze (the earliest are at 00:08 and
00:23~UTC), and the statement files were not changed afterwards.

\paragraph{The proofs.}
For Theorems~\ref{thm:two} and~\ref{thm:five} the parameter values are
computed with \texttt{decide~+kernel}, the rank of the five-point group is
bounded below by showing that it is not abelian, and the frozen statements are
proved by \texttt{TaskJ.\allowbreak Rank.\allowbreak Proof.\allowbreak
TwoPointCounterexample} and
\texttt{TaskJ.\allowbreak Rank.\allowbreak Proof.\allowbreak
FivePointCounterexample} (and four further theorems with the names of the
other targets). For Theorem~\ref{thm:gap}, the labels are
$p_{ij}={}$\texttt{Nat.nth Nat.Prime (Nat.pair i j)} for $i\ne j$, and
\texttt{TaskJ.\allowbreak Rank.\allowbreak Extension.\allowbreak Proof.\allowbreak
UnboundedGap} combines the arithmetic of the labels
(\texttt{Arithmetic.lean}), the orbits and stabilizers
(\texttt{OrbitStructure.lean}), faithfulness by the Chinese remainder theorem
(\texttt{RankCRT.lean}), the identification of the equivariant group with the
product of the translation groups (\texttt{FiberAutomorphisms.lean}), the rank
of $C_M\times C_M$ (\texttt{SquareRank.lean}) and the packet count for every
valid choice (\texttt{PacketRank.lean}). The test files
\texttt{CheckRank.lean} and \texttt{CheckExtension.lean} state each frozen
target with its frozen type, as \texttt{theorem check\_X : Challenge.X :=
Proof.X} (for \texttt{UnboundedGap} also under its full frozen name). The axioms of every proof and check theorem, printed by the files,
are exactly \texttt{propext}, \texttt{Classical.choice} and
\texttt{Quot.sound}. No file uses \texttt{sorry}, \texttt{native\_\allowbreak
decide}, \texttt{bv\_\allowbreak decide} or an \texttt{axiom} declaration;
\texttt{decide~+kernel} is checked by the kernel and adds no axiom.

\paragraph{Checks.}
Each of the two packages was compiled module by module, from its sources, in
two new empty output directories on the author's Mac, with the pinned
toolchain and Mathlib; every module was replayed with \texttt{leanchecker} in
a separate process, and each directory also built the whole package. Mathlib
and its dependencies were taken from a shared prebuilt cache at the pinned
commits and were not rechecked; the \texttt{leanchecker} replay covers the
$15$ project modules. The package for Theorems~\ref{thm:two}
and~\ref{thm:five} ($5$ modules) was checked on 18 November 2026 from 00:14 to
00:16~UTC with $22$ commands, and the package for Theorem~\ref{thm:gap}
($10$ modules) from 00:29 to 00:32~UTC with $42$ commands. All $64$ commands
ended with exit code~$0$, all \texttt{leanchecker} outputs were empty, the
builds printed nothing except the requested axiom lists, and the compiled
modules of the two directories have the same SHA-256 hashes. (Of two earlier runs of the
first package, at 00:10 and 00:11~UTC, the first stopped before compiling
because the Mathlib cache was not found and the second passed.)
In the repository the files have module names under \texttt{Research.TaskJ} and differ from the checked files only in their \texttt{import} lines; these versions were compiled and replayed with \texttt{leanchecker} again, module by module, in two new empty directories on the author's Mac (18 November 2026, 19:12--19:21~UTC), with identical \texttt{.olean} files and the same axiom lists. In the whole-project rebuild of version~1.23.0 of the repository (18--19 November 2026, on the author's Windows computer), these files were compiled and replayed with \texttt{leanchecker} again, together with the rest of the project, each at the first attempt (details in the repository). The formalization has not been checked with comparator or with
an independent kernel.

\section*{How this paper was produced}
This work was carried out from 17 to 19 November 2026 (UTC) in a workflow
directed by the author, using OpenAI Codex with the model GPT-6.2-Astra
(reasoning level ultra) and Anthropic Claude Code with the model Claude
Opus~5.5. The author asked for a search for open problems with the newest
Codex model, and Claude Code passed this request to Codex. Between 19:34 and
20:23~UTC on 17 November, Codex screened $45$ problems and sub-questions from
the recent literature, wrote proof sketches for three of them and proposed
them; the author chose this one and one other. From 23:49~UTC on 17 November
Codex agents with the same model did the following. One agent, which wrote
neither the statement files nor the proofs, checked the reading of both
versions of the source and then reviewed the two statement files before they
were frozen; these three checks were thus done by a single agent and are not
independent of each other. Other agents wrote the statement files, the written
proofs and the Lean proofs, and ran the checks, which ended at 00:32~UTC on
18 November. The results were then reviewed by an independent Claude Code
agent (model Claude Opus~5.5) that had written none of the files. It read the
preprint itself and the journal version in an archived copy (it could not
open the publisher's page), read the statement files, the check files and the main
proof files, read the receipts and the raw logs of the checks, and recomputed
the parameters of the examples with its own programs; it did not read the files \texttt{OrbitStructure.lean} and
\texttt{FiberAutomorphisms.lean} line by line. On 18 November at 01:10~UTC it
approved the formal results (APPROVE\_FORMAL) with five minor remarks on
wording and attribution, which were applied to Codex's written proofs and to
this text. Claude Code (model Claude Opus~5.5) checked the passages of the
sources cited here (those of the journal version in the same archived copy)
and wrote this text from Codex's manuscript and the Lean
files. An independent Claude Code agent (model Claude Opus~5.5) that had written none of the files reviewed this paper, and another reviewed the release; their corrections were applied. Throughout, ``we''
refers to this workflow; the reading and searching described in the paper
were done by the agents. The author is not a professional mathematician.
The author has read this paper and, using a side-by-side table, compared its Lean statements with Theorems~\ref{thm:two}, \ref{thm:five} and~\ref{thm:gap}. No human mathematician has reviewed this paper. The correctness
of Theorems~\ref{thm:two}, \ref{thm:five} and~\ref{thm:gap} rests on the Lean
proofs and on the agreement of the Lean statements with the definitions
of~\cite{lrr,rlr}; the written proofs are given for the reader, and the
remarks marked as written arguments are not formally verified. AI systems are
not authors of this paper; the author takes full responsibility for its
content.

\medskip
\noindent Thanks are due to Victor M. Lara-G\'omez, Gerardo Romero-Rosales and
Ram\'on H. Ruiz-Medina for their paper and its explicit constructions, on
which this note relies throughout.

\begin{thebibliography}{9}
\bibitem{repo} H.~Dong, \emph{ai4math-results}, repository,
\url{https://github.com/Horace-Maxwell/ai4math-results}, version~\repoversion{}
(2026); archived at Zenodo under the concept DOI
\url{https://doi.org/10.5281/zenodo.22970254}.
\bibitem{lrr} V.~M. Lara-G\'omez, G.~Romero-Rosales and R.~H. Ruiz-Medina,
Finite generating sets for monoids of $G$-equivariant functions,
\emph{Semigroup Forum}, published online 2 November 2026,
\url{https://doi.org/10.1007/s00233-026-10690-y}.
\bibitem{lu} J.~Lu, Generation of generalised wreath products of symmetric
groups, arXiv:2604.16141v1, 2026.
\bibitem{mathlib}
The mathlib Community, The Lean mathematical library, in \emph{Proc.\ 9th ACM
SIGPLAN Int.\ Conf.\ Certified Programs and Proofs} (CPP '20), ACM, 2020,
\url{https://doi.org/10.1145/3372885.3373824}.
\bibitem{rlr} R.~H. Ruiz-Medina, V.~M. Lara-G\'omez and G.~Romero-Rosales,
Finite generating sets for monoids of $G$-equivariant functions,
arXiv:2607.04053v1, 2026.
\end{thebibliography}
\end{document}
